\section{Setting and definitions}
\label{sec:setting}
\paragraph{Leaderboard and rules.}
A leaderboard is a score matrix $X\in[0,100]^{M\times n}$ of $M$ models on $n$ tasks, taken from a fixed public snapshot. An aggregation rule $f$ induces a ranking; we study the \emph{score rules} arithmetic mean, median (mean of the middle pair for even $n$), minimum, and shifted geometric and harmonic means (computed on score$+1$ so that zeros are admissible, then shifted back), and the \emph{ordinal rules} Borda (summed per-task ranks), plurality (per-task first places, ties split equally), and, as a requirement rather than a ranking, the \emph{Condorcet requirement} (beat every rival in a strict task majority, task ties counting half). Because a Condorcet-consistent rule can also elect a model that is not a Condorcet winner, the Condorcet cost is an upper bound for such rules, not their cost. Trimmed-mean and power-mean variants are sensitivity checks.

\paragraph{Investment.}
Fix a target model with score vector $x\in\mathbb{R}^n$ and hold all rival rows fixed. Targeted investment raises the target's score on task $t$ to $y_t\in[x_t,c_t]$, where $c_t$ is a \emph{cap} (the effective cap is $\max(x_t,\text{scenario cap})$, with scenario caps clipped at 100, so a target that already exceeds a scenario cap on a task keeps its original score there instead of facing an empty box); this stands for training on data related to sub-dataset $t$. The target is \emph{first} if $f(y)\ge \max_{r}f(x_r)+\varepsilon$ for score rules ($\varepsilon=0.001$ points; ordinal rules use an analogous finite convention: per-task levels are the original score, a tie with a rival score, or one margin step above it, taking only rival scores at or above the target's original score on that task), and a target that is already uniquely first has cost zero. We measure investment in three ways, each optimized separately: $K=\|y-x\|_0$ (tasks touched), $L_1=\sum_t(y_t-x_t)$ (score points) and $L_\infty=\max_t(y_t-x_t)$ (the largest single-task jump, which is the budget of a uniform ``general improvement''). Their minimizers need not coincide. Infeasible cases (no vector within the caps makes the target first) have infinite cost and are reported as a feasible fraction, never dropped.

\paragraph{Caps are the modelling assumption.}
Public scores identify no achievable training gain. We therefore treat the cap as a scenario and show every result with it: \texttt{ceil} ($c_t=100$; what is needed regardless of realism), \texttt{front0/2} (the best rival score on $t$ plus 0 or 2 points), \texttt{gain5/10/20} (at most $+5/10/20$ points per task) and, for the Open LLM Leaderboard only, \texttt{emp95/99}, per-task quantiles of observed fine-tune-minus-base score differences. The latter are descriptive scales, confounded by selection, architecture and training, and are \emph{not} estimates of achievable gains.

\paragraph{Exact results.}
The following hold for the model above (proofs in Appendix~\ref{app:proofs}; \texttt{analysis/} checks each against the solvers). Throughout, a target that is already uniquely first costs 0 by convention, even if it wins by less than $\varepsilon$; the formulas below are for the other targets. (i) \emph{Mean:} with gap $g=\tau+\varepsilon-\bar x$ to the best rival mean $\tau$, minimum $L_1=n\,g$ whenever total headroom suffices, and minimum $K$ is the shortest descending-headroom prefix reaching $n\,g$. (ii) \emph{Minimum:} with $b=$ best rival minimum $+\varepsilon$, feasibility requires every cap $\ge b$; then $K=\#\{t:x_t<b\}$ and $L_1=\sum_t\max(0,b-x_t)$. (iii) \emph{Odd median} ($n=2h-1$): at least $h$ final coordinates must reach $b$; $L_1$ is the sum of the $h$ cheapest \emph{cap-feasible} lifts (infeasible if fewer than $h$ coordinates can reach $b$). (iv) \emph{Cap threshold:} for any rule whose set of first-placing plans in the ceiling box is closed (all rules studied here), a gain-$h$ scenario is feasible iff the $\texttt{ceil}$ value of $L_\infty$ is $\le h$. (v) \emph{No universal order:} the costs of the mean, median and minimum rules can reverse (Table~\ref{tab:counter}), so no universal ordering of these rules and no blanket statement such as ``the median needs half the tasks'' holds; we do not prove reversals for every other rule studied. (vi) \emph{Direction of the bounds:} the score-rule minima are exact minima in this fixed-rival box model, hence lower bounds on the score-space cost of any intervention feasible in the same box, not on the cost of an actual training plan. The ordinal minima are exact only on the declared finite set of per-task levels (unchanged, tie a rival, pass a rival by one margin step), so they bound the infimum over arbitrary real scores from above, and infeasibility on that set is conditional on it. The Condorcet requirement gives a different, sufficient-condition upper bound for a Condorcet-consistent rule.
\begin{table}[t]
\caption{Constructed (not measured) examples with $n=5$ tasks, one rival row and margin $\varepsilon=10^{-3}$: exact minimum $K$ and $L_1$ for the target to be first. The cheapest rule changes between examples; in the last one the target already wins under the median (cost 0). Source: \texttt{analysis/results/theory\_examples.json}.}
\label{tab:counter}
\centering\scriptsize\setlength{\tabcolsep}{3pt}
\input{../tables/tab_counterexamples.tex}
\end{table}
